\documentclass[10pt,a4paper]{article}

\usepackage[margin=1.05in]{geometry}
\usepackage{setspace}
\usepackage{amsmath,amssymb,amsthm}
\usepackage{graphicx}
\usepackage{enumitem}
\usepackage{titlesec}
\usepackage[expansion=false]{microtype}
\usepackage{fancyhdr}
\usepackage{hyperref}
\hypersetup{colorlinks=true,linkcolor=black,urlcolor=black,citecolor=black}

\graphicspath{{../../drawings/}{../../drawings/sheets/}}

\titleformat{\section}{\normalsize\bfseries}{}{0em}{}
\titleformat{\subsection}{\small\bfseries}{}{0em}{}
\titlespacing*{\section}{0pt}{1.35em}{0.5em}
\titlespacing*{\subsection}{0pt}{0.9em}{0.3em}

\theoremstyle{plain}
\newtheorem{proposition}{Proposition}
\theoremstyle{definition}
\newtheorem{definition}{Definition}
\theoremstyle{remark}
\newtheorem{observation}{Observation}

\setstretch{1.14}
\setlist[itemize]{leftmargin=1.15em,itemsep=0.1em,topsep=0.2em}

\pagestyle{fancy}
\fancyhf{}
\fancyhead[L]{\small\itshape Elements of Position}
\fancyhead[R]{\small\itshape Without a Landing}
\fancyfoot[C]{\small\thepage}
\renewcommand{\headrulewidth}{0pt}

\begin{document}

\begin{center}
\vspace*{1.6em}
{\large Elements of Position}\\[1.4em]
{\LARGE\bfseries Without a Landing}\\[0.5em]
{\small The residual the operators leave, the pair they still write, and which meetings hold}\\[1.5em]
{\normalsize Justin Erholtz}\\[0.3em]
{\small 19 September 2026}
\end{center}

\vspace{1.2em}

\begin{center}
{\large The operators do not contain the joint.}
\end{center}

\vspace{1.2em}

\noindent
\textit{A note.}
Part~I named an entrance residual and then, in places, treated it as dirt on an idealization.
This part writes the residual as an object, and then stays with the pair long enough to see what the pair still forces after the landing is refused.
Claims of two kinds stay apart.
A heading identity is checked against \(u_m\).
A pair identity is checked against \(m\).
A count is reported from a table.
A length that appears only after a table exists may be named as a comparison, and not before.
No angle enters the step.

\vspace{0.5em}
\noindent
\textit{A note on names used here.}
The names of Part~I stand: shear, bind, step, hit, pair, lock, leftover, inversion.
Added below, only as readings of those objects: mouth of the kite; inner and outer length; kite factor \(1-1/m^2\); new inner --- an inner length that is not a product of earlier inner lengths; first split --- the hit \(m=2\).
The recorded run cited throughout is still \(h=0.002\), \(200\,000\) steps, \(127\) crossings, unless a different kick is named.

\tableofcontents

% ======================================================================
\section{The leftover of the landing is forced}
% ======================================================================

Hit \(m\) occurs at step \(j_m=\lceil m\lambda\rceil\),
\(\lambda=\pi/\arctan h\), whenever \(m\lambda\) is not an integer.
The heading has already passed \(m\pi\) by
\[
\varepsilon_m=\arctan(h)\,(1-\{m\lambda\}).
\]

\begin{proposition}
\[
u_m=\bigl((-1)^m\cos\varepsilon_m,\;(-1)^m\sin\varepsilon_m\bigr).
\]
\end{proposition}
Checked on \(827\) hits at \(h=0.002\), residual \(3\times10^{-13}\);
on \(397\) hits at \(h=0.005\), residual \(10^{-13}\).
The formula does not use \(0.002\).

\begin{proposition}
\(\varepsilon_m=0\) if and only if \(m\lambda\in\mathbb{Z}\).
\end{proposition}
That is the exact-closure kick \(h=\tan(\pi/n)\) of Part~I.
At \(n=100\), \(2000\) hits, one gap, \(\max\varepsilon=9\times10^{-13}\).

The operators do not contain the joint.
They also do not contain the wall.
A hit is the first sample after the crossing.

\begin{figure}[htbp]
\centering
\includegraphics[width=0.92\textwidth]{E1R_residual_elevation.png}
\caption{Sheet E1R. Elevation of the residual. The wrap and the defect schedule sit on the same walk that wrote the pair.}
\end{figure}

% ======================================================================
\section{The schedule is a continued fraction}
% ======================================================================

After the first hit, the gap is short if and only if
\[
\{(m-1)\lambda\}<1-\{\lambda\}.
\]
\(826\) hits, no failures.

On the recorded kick, \(1-\{\lambda\}=[0;4,1,24,\ldots]\).
Spacings \(4\) and \(5\). Defects alternate on the next denominators
\(124\) and \(129\).
The number \(\alpha\) of Part~I is the slope of the first block
\(m=5,10,\ldots,125\) only.
The first defect kills it.
Each later block is a new ramp.
Those are readings of one kick's continued fraction.
A different kick writes a different word.
It does not revise the heading formula.

% ======================================================================
\section{What does not need a landing}
% ======================================================================

These are identities of the pair and of inversion.
They use \(u_m\) only as a ray.
They do not care that \(u_m\) missed the cut.

\begin{align*}
|OP_m|\,|OQ_m|&=1,\\
Q_m&=P_m/|P_m|^2,\\
\bigl(\tfrac1m,\,m;\,1,\,-1\bigr)&=-1,\\
\bigl(m-2,\,m+2;\,m-1,\,m+1\bigr)&=\tfrac19,\\
\tau(m)&=\tfrac{(m-1)^2}{2m}.
\end{align*}

\begin{proposition}
The circle on diameter \(OP_m\) meets the unit circle at points
whose foot on the ray of \(u_m\) is exactly \(Q_m\).
\end{proposition}
Checked on the recorded headings, foot error \(2\times10^{-16}\).
No idealization. The lock is the case \(m=1\), where the two circles
are one circle and the foot is the point itself.
Same construction on a lock, a pair, a long hit, and a short hit:
the foot is \(Q\) on every residual.

The two meeting-points on the bind are \(X_\pm\).
The segment \(X_+X_-\) is the polar of \(P_m\).
The angle at each \(X\) is a right angle.
Checked on all \(127\) recorded headings: the departure from a right angle is \(0\) to machine precision, lock included.

\begin{figure}[htbp]
\centering
\includegraphics[width=0.92\textwidth]{D6R_kite_dirty.png}
\caption{Sheet D6R. The kite on the recorded headings. The residual is visible. The foot is still \(Q\).}
\end{figure}

\begin{figure}[htbp]
\centering
\includegraphics[width=0.92\textwidth]{K1_kite_stacked.png}
\caption{Sheet K1. The same kite, aligned and dirty, stacked so the construction can be read at once.}
\end{figure}

\begin{figure}[htbp]
\centering
\includegraphics[width=0.92\textwidth]{O2_kite_solid.png}
\caption{Sheet O2. The kite read as a solid: outer cone, bind as cylinder, inner funnel, polar as the inverted circle on \(OP\).}
\end{figure}

\begin{figure}[htbp]
\centering
\includegraphics[width=0.92\textwidth]{D8_polar.png}
\caption{Sheet D8. The circle on \(OP\), inverted through the bind, is the polar line \(X_+X_-\).}
\end{figure}

The mouth of the kite --- the length \(X_+X_-\) --- is
\[
2\sqrt{1-\frac1{m^2}}.
\]
At the lock the mouth is \(0\).
At the first split it is \(\sqrt{3}\).
It walks toward \(2\) and does not reach \(2\) at any finite hit.

% ======================================================================
\section{What needs a landing}
% ======================================================================

Opposite-wall chords
\[
|P_m-P_{m-1}|=2m-1,\qquad
|Q_m-Q_{m-1}|=\frac{2m-1}{m(m-1)}
\]
assume \(u_m=-u_{m-1}\).
The operators deliver
\[
u_m\cdot u_{m-1}=-\cos(\varepsilon_m-\varepsilon_{m-1}).
\]
The missing length is of order \(\varepsilon^2\).
On a long gap of the recorded run the relative miss is \(2\times10^{-8}\).
On a short gap it is larger, order \(10^{-7}\) to \(10^{-5}\).
The product of the two actual chords approaches \(4\) and is not \(4\)
at any finite hit: at \(m=127\), the miss from \(4\) is \(6\times10^{-5}\).

The swap of a collinear diameter-circle assumes two headings
parallel or antiparallel.
Same-wall pairs of the recorded run have \(1-|u_m\cdot u_n|\)
between \(10^{-11}\) and \(10^{-7}\), set by how close those two
residuals are to each other, not by how far apart the hits are.

% ======================================================================
\section{The mix family}
% ======================================================================

\(T_{a,b}(x,y)=(x-ay,\,y+bx)\), then bind.
Opposite signs return, two gaps, and
\(\lambda=\pi/\arctan\sqrt{ab}\).
Equal \(a=b\) is balance, not a new kind.
Same-sign Euclidean bind does not return.
One-sided mix does not return.
Hyperbolic bind to \(x^2-y^2=1\) does not return.

The minus is the kind.
The residual law rides on whatever kick still returns.
The equal opposite shear is a right angle in the mix.
It is not a finished turn in the mix.
The right angle at \(X\) on the bind is the other place that angle appears.
The walk never puts it in the step.

\begin{figure}[htbp]
\centering
\includegraphics[width=0.92\textwidth]{SCH3_stands.png}
\caption{Sheet SCH3. What survives a change of mix, and what was only a reading of one kick.}
\end{figure}

% ======================================================================
\section{Three readings of one pair}
% ======================================================================

Once \(m\) is named, three numbers sit on the pair and do not ask the heading anything further.

\begin{itemize}
\item Inner length \(|Q_m|=1/m\).
\item Leftover \(\tau(m)=(m-1)^2/(2m)\).
\item Mouth \(2\sqrt{1-1/m^2}\).
\end{itemize}

Lock: leftover \(0\), mouth \(0\), inner \(=\) outer \(=1\).
First split, \(m=2\): leftover \(1/4\), mouth \(\sqrt{3}\), inner \(1/2\).
After that they are different jobs.
Inner falls toward the cut.
Mouth is already near \(2\) by \(m=10\).
Leftover keeps climbing, like \(m/2\).

The right angle at \(X\) is not among them.
It would be a flat line.
That is the distinction Part~I already wanted: the bind holds a right angle; the pair holds how far the hit is from lock.

\begin{figure}[htbp]
\centering
\includegraphics[width=0.92\textwidth]{P1_pair_leftover_mouth.png}
\caption{Sheet P1. Leftover grows. Mouth jumps at the first split and then sits near \(2\). The right angle on the bind does not appear because it does not move.}
\end{figure}

\begin{figure}[htbp]
\centering
\includegraphics[width=0.92\textwidth]{P2_three_readings.png}
\caption{Sheet P2. Inner, leftover, and mouth on one hit. They sit together only at the lock.}
\end{figure}

% ======================================================================
\section{Inversion does not move}
% ======================================================================

\begin{proposition}
For every hit \(m\ge 1\),
\[
|Q_m|\,|P_m|=1.
\]
\end{proposition}
The product is the pair.
Outer climbs. Inner falls. Their product is \(1\) at the lock and \(1\) after the lock.
Leftover is not a factor in that product.
Mouth is not either.

Two constants from one hit, then:
the product of the two lengths, and the right angle on the bind.
One lives on the pair.
One lives on the bind.
Neither was put in the shear.

\begin{figure}[htbp]
\centering
\includegraphics[width=0.92\textwidth]{P3_inversion.png}
\caption{Sheet P3. Inner times outer is \(1\) at every hit.}
\end{figure}

% ======================================================================
\section{A new inner is not a product of earlier inners}
% ======================================================================

The lengths multiply.
\[
|P_{ab}|=|P_a|\,|P_b|,\qquad |Q_{ab}|=|Q_a|\,|Q_b|.
\]
Checked through the recorded table, error \(0\).
So \(6\) is the product of the pairs at \(2\) and at \(3\).
\(4\) is the pair at \(2\), taken twice.
\(15\) is the pairs at \(3\) and at \(5\).

A hit whose inner length is not a product of earlier inner lengths is a new inner.
The first of those is the first split, \(m=2\).
Then \(3,5,7,11,\ldots\)
On the recorded run: one lock, thirty-one new inners, twelve powers of a single new inner, and the rest mixed.

Leftover does not see this cut.
Every \(m>1\) has leftover.
The leftover curve of Part~I is smooth through new inners and through products alike.
Thales does not see it either.
Foot error and polar error stay at \(10^{-16}\) on lock, new inner, power, and mixed.

Two geometries share the list of hits and ignore each other.
The order of hits writes the short-long word.
The lengths of hits write which inners are new.
A cross of the two, on several kicks, sits on independence except where a given kick's word happens to live on a multiple --- the recorded five-block, for instance.
That is a schedule accident, not a second law.

\begin{figure}[htbp]
\centering
\includegraphics[width=0.92\textwidth]{F1_irreducible_pairs.png}
\caption{Sheet F1. Height is how many new-inner factors make the length. Accent is a new inner. The leftover panel underneath does not read that height.}
\end{figure}

\begin{figure}[htbp]
\centering
\includegraphics[width=0.92\textwidth]{F3_thales_kinds.png}
\caption{Sheet F3. The same kite at lock, a new inner, a power, and a mixed hit. Thales does not read the kind.}
\end{figure}

% ======================================================================
\section{The kite factor}
% ======================================================================

At every hit the pair writes
\[
1-\frac1{m^2}.
\]
That is the kite factor: the square of half the mouth, and also \(1-|Q_m|^2\).

The product of those factors from \(m=2\) to \(N\) telescopes:
\[
\prod_{m=2}^{N}\Bigl(1-\frac1{m^2}\Bigr)=\frac{N+1}{2N}.
\]
It walks toward \(1/2\).
Integers only.
No named constant is required to see the floor.

The same factor, kept only at new inners, walks toward another floor.
At \(N=80\) the running product is \(0.60926\) against the named comparison \(6/\pi^2=0.60793\).
The name is a sheet reading.
The object is the product of kite factors at hits whose inner is not a product.

The same factor, kept only at the squares of new inners --- \(4,9,25,49,\ldots\) --- is \(1-1/p^4\) when the new inner is \(1/p\).
Those factors sit near \(1\).
Their product walks toward the named comparison \(90/\pi^4\).
Powers as a class are sparse.
They barely pull.

\begin{figure}[htbp]
\centering
\includegraphics[width=0.92\textwidth]{P4_new_inners.png}
\caption{Sheet P4. Accent marks a new inner. Below: product of kite factors at every hit, falling to \(1/2\); product at new inners only, falling toward \(6/\pi^2\).}
\end{figure}

\begin{figure}[htbp]
\centering
\includegraphics[width=0.92\textwidth]{P5_powers.png}
\caption{Sheet P5. Squares mark a power of one new inner. Their kite factors stay near \(1\). The third path barely moves.}
\end{figure}

% ======================================================================
\section{The telescope factors}
% ======================================================================

Write \(A_N\) for the product of kite factors at every hit from \(2\) to \(N\),
\(N_N\) for the same product kept only at new inners,
and \(C_N\) for the product kept at every remaining hit --- powers and mixed together.

\begin{proposition}
\[
A_N=N_N\cdot C_N
\]
at every \(N\ge 2\).
\end{proposition}
Checked to \(N=400\), residual \(10^{-16}\).
The identity is the partition of the same list, not a new operator.

The floors those three products walk toward stand in the same relation:
\[
\frac12=\frac{6}{\pi^2}\cdot\frac{\pi^2}{12}.
\]
The names cancel.
What remains is the telescope, factored through the cut ``is this inner new, or not.''

Not-new is powers plus mixed.
Powers barely pull.
Mixed does most of the not-new work.
The first mixed hit is \(m=6\), inner \((1/2)(1/3)\), kite factor \(35/36\).

The number \(\pi^2/12\) is also the named comparison for the wall-sign sum of inner squares --- opposite hits, alternating.
Two writings, one number, set next to each other after both exist.
They are not declared to be one construction.

\begin{figure}[htbp]
\centering
\includegraphics[width=0.92\textwidth]{P6_telescope_factors.png}
\caption{Sheet P6. All hits, new inners, and not-new. One kite factor. Three floors. The identity \(A_N=N_N\cdot C_N\) holds at every finite \(N\).}
\end{figure}

% ======================================================================
\section{Exact closure, nested}
% ======================================================================

When the kick is \(h=\tan(\pi/n)\) the walk lands.
One gap.
The period in hits is \(n\).
Headings of a smaller exact closure sit inside a larger one whenever the smaller period divides the larger.

Checked: every heading of the \(n=3\) walk sits on the \(n=6\) walk at even steps, chord \(10^{-16}\).
Every heading of \(n=4\) sits on \(n=8\) the same way.
The square \(h=1\) is the case \(n=4\).
It is the only rational kick that lands among those tried in Part~I.

A kick cannot be written for \(n=2\): that half-turn is the pole of the tan-addition that doubles a period.
Finite large \(h\) approaches a right-angle step and never arrives.
Two steps approach the far side of the cut and stay a hair above it, so the first recorded hit of a large finite kick remains the third step.
The kite does not need that kick.
It already writes the right angle at \(X\) on every finite hit.

\begin{figure}[htbp]
\centering
\includegraphics[width=0.92\textwidth]{F4_nested_closures.png}
\caption{Sheet F4. Exact closures nested: \(3\subset 6\subset 12\). The square \(n=4\) is the only rational landing.}
\end{figure}

% ======================================================================
\section{How often a new inner appears}
% ======================================================================

Among the first \(N\) hits, how many inners are new?
At \(N=127\), thirty-one.
At \(N=400\), seventy-eight.
The fraction falls.
A named comparison for that count is \(N/\log N\).
At \(N=400\) the ratio of the count to that comparison is about \(1.17\).
The name is a sheet reading of how often a new inner appears among the magnitudes the walk already wrote.
It is not put in the shear.

\begin{figure}[htbp]
\centering
\includegraphics[width=0.92\textwidth]{F2_pair_monoid.png}
\caption{Sheet F2. The running count of new inners, and the four kinds of hit on the recorded run: lock, new inner, power of one new inner, mixed.}
\end{figure}

% ======================================================================
\section{What stands}
% ======================================================================

Forced at every hit, residual included:
the heading formula;
the pair, the leftover, the inversion, the kite-foot, the right angle at \(X\);
the short-gap rule as a fractional part;
exact closure exactly when \(m\lambda\) is an integer;
the product of inner and outer;
the kite factor and the telescope of all kite factors;
the identity that the telescope is the product of the new-inner product with the rest.

Forced only as well as two headings agree:
opposite-wall chords, their product, the swap.

Reported from a table:
the first-block slope \(\alpha\), the defect list, the wrap;
the running products at finite \(N\), and the named floors they walk toward;
the count of new inners among \(N\) hits.

Named: entrance residual, defect, wrap, first block, mouth, kite factor, new inner, first split.

Not granted:
a landing on the cut;
\(\alpha\) as a constant of the operators;
a value of \(4\), or of \(1/2\), or of \(2\), treated as reached at a finite hit;
a name treated as a step.

The first sentence is still the sentence.
Where is \(1\)?
It is the lock.
The residual at the lock is just another \(\varepsilon_1\).
Even there the operators did not land.
They only met themselves.
Inner and outer are one point.
The leftover is \(0\).
The mouth is \(0\).
The right angle on the bind is already there, and has nowhere to open.
Everything after is the same pair, walked out.

\end{document}
